\documentclass[12pt,a4paper]{article}
\usepackage[margin=22mm]{geometry}
\usepackage{fontspec}
\setmainfont{FreeSerif}
\newfontfamily\thaifont{FreeSerif}
\usepackage{amsmath,amssymb,tikz,enumitem}
\usetikzlibrary{arrows.meta,patterns}
\usepackage{fancyhdr}
\pagestyle{fancy}\fancyhf{}\fancyhead[L]{Physics Olympiad — Problem Bank}\fancyhead[R]{21–30}\fancyfoot[C]{\thepage}
\XeTeXlinebreaklocale "th"
\XeTeXlinebreakskip=0pt plus 1pt
\setlength{\parindent}{0pt}\setlength{\parskip}{7pt}
\begin{document}
\begin{center}{\Large โจทย์ฟิสิกส์ ข้อ 21–30}\\[3pt] {\small ฉบับถอดใหม่พร้อมแผนภาพ TikZ}\end{center}
\section*{ข้อ 21 \hfill \normalsize MEC-0021}
(สอวน ม.4 2557) ระยะทาง $AB=D$ เป็นรางรถไฟตรง นกออกบินจาก $A$ พร้อมกันกับที่รถไฟออกจาก $B$ นกบินเร็ว $v$ รถไฟเคลื่อนที่เร็ว $u$ เมื่อนกถึงรถไฟครั้งแรกก็บินกลับทันทีด้วยความเร็วขนาดเท่าเดิม และเมื่อถึงจุด $A$ นกก็หันกลับบินเข้าหารถไฟอีก และเมื่อถึงรถไฟก็บินกลับอีก กลับไปกลับมาแบบนี้จนกระทั่งรถไฟถึงจุด $A$ รวมระยะทางทั้งหมดที่นกบินได้เป็นเท่าใด (กำหนดว่า $v>u$)
\begin{center}\begin{tikzpicture}[>=Stealth,scale=.8]\draw[dashed] (-.5,0)--(8,0);\draw[thick] (1,0)--(7,0);\fill (1,0) circle(2pt) node[below]{$A$};\fill (7,0) circle(2pt) node[below]{$B$};\node[above] at (1.7,.12){นก};\node[above] at (6.5,.12){รถไฟ};\draw[->,thick] (1.5,.55)--(3.1,.55) node[midway,above]{$v$};\draw[->,thick] (6.8,.55)--(5.2,.55) node[midway,above]{$u$};\end{tikzpicture}\end{center}
\medskip

\section*{ข้อ 22 \hfill \normalsize MEC-0022}
(สอวน ม.5 2557) ระยะทาง $AB=D$ เป็นรางรถไฟตรง นกออกบินจาก $A$ พร้อมกันกับที่รถไฟออกจาก $B$ นกบินได้เร็ว $v$ รถไฟเคลื่อนที่เร็ว $u$ เมื่อนกถึงรถไฟครั้งแรกก็บินกลับทันทีด้วยความเร็วขนาดเท่าเดิม เมื่อถึงจุด $A$ รถไฟจะอยู่ห่าง $B$ เท่าใด (กำหนดว่า $v>u$ นกจะได้ไม่ถูกชน)
\begin{center}\begin{tikzpicture}[>=Stealth,scale=.8]\draw[dashed] (-.5,0)--(8,0);\draw[thick] (1,0)--(7,0);\fill (1,0) circle(2pt) node[below]{$A$};\fill (7,0) circle(2pt) node[below]{$B$};\node[above] at (1.7,.12){นก};\node[above] at (6.5,.12){รถไฟ};\draw[->,thick] (1.5,.55)--(3.1,.55) node[midway,above]{$v$};\draw[->,thick] (6.8,.55)--(5.2,.55) node[midway,above]{$u$};\end{tikzpicture}\end{center}
\medskip

\section*{ข้อ 23 \hfill \normalsize MEC-0023}
(สอวน ม.4 2552) มีรถเมล์ออกจากบ้านจันสองหน้าไปยังบ้านสี่สาวเป็นประจำ หนึ่งคันทุกๆ 10 นาที บ้านสี่สาวอยู่ห่างจากบ้านจันสองหน้าเป็นระยะทาง $60\ \mathrm{km}$ รถเมล์แล่นด้วยอัตราเร็วคงตัว $60\ \mathrm{km/h}$ วันหนึ่งนายเอกขับรถด้วยอัตราเร็ว $60\ \mathrm{km/h}$ ออกจากบ้านสี่สาวไปบ้านจันสองหน้าเวลาเดียวกับที่รถเมล์คันหนึ่งกำลังออกจากบ้านจันสองหน้าพอดี เขาจะสวนกับรถเมล์กี่คันระหว่างการเดินทาง (ไม่นับตอนตั้งต้นและที่ปลายทาง)
\medskip

\section*{ข้อ 24 \hfill \normalsize MEC-0024}
(สอวน ม.4 2548) จอมตลกวางแผนปล่อยหนูปีศาจออกมาอาละวาดในเมือง เขาเอาหนูปีศาจใส่รถบรรทุกแล่นด้วยความเร็ว $10\ \mathrm{m/s}$ ไปตามถนนตรง และปล่อยหนูปีศาจออกมาทีละตัว โดยปล่อยหนึ่งตัวต่อวินาที หนูปีศาจเมื่อหลุดจากรถก็วิ่งเข้าไปข้างหน้าทันทีด้วยอัตราเร็ว $20\ \mathrm{m/s}$ (อัตราเร็วของหนูปีศาจบนถนนไม่ขึ้นกับอัตราเร็วของรถบรรทุก) มนุษย์ค้างคาวรู้ข่าวนี้ล่วงหน้า จึงรีบมาจัดการเก็บหนูปีศาจ เขาขับรถค้างคาวแล่นสวนเข้ามาด้วยความเร็ว $20\ \mathrm{m/s}$ ในทิศตรงข้ามกับทิศที่หนูวิ่ง รถค้างคาวมีเครื่องมือเก็บหนูปีศาจทันทีที่หนูถึงรถ จงหาว่ามนุษย์ค้างคาวเก็บหนูปีศาจไปทั้งหมดเฉลี่ยได้กี่ตัวต่อหนึ่งนาที
\medskip

\section*{ข้อ 25 \hfill \normalsize MEC-0025}
(สอวน ม.4 2554) วัตถุเคลื่อนที่ด้วยอัตราเร็วคงที่ $v_1$ จากจุด $A$ เป็นระยะทาง $a$ แล้วเปลี่ยนอัตราเร็วเป็น $v_2$ และเคลื่อนที่ต่อไปเป็นระยะทางอีก $b$ จนถึงจุด $B$ อัตราเร็วเฉลี่ยของวัตถุในการเคลื่อนที่จาก $A$ ไป $B$ เป็นเท่าใด
\medskip

\section*{ข้อ 26 \hfill \normalsize MEC-0026}
(IJSO รอบที่ 1 2555) วัตถุชิ้นหนึ่งเคลื่อนที่ในหนึ่งมิติ จากจุด $A$ ไปยังจุด $B$ ด้วยอัตราเร็วคงตัว $6.0\ \mathrm{m/s}$ แล้วเคลื่อนที่กลับมาที่จุด $A$ ด้วยอัตราเร็วคงตัว $4.0\ \mathrm{m/s}$ ความเร็วเฉลี่ยของวัตถุนี้ในช่วงเวลาตั้งแต่เคลื่อนที่ออกจากจุด $A$ แล้วกลับมาที่จุด $A$ อีกครั้งหนึ่งเป็นเท่าใด
\begin{enumerate}[label=\arabic*.,itemsep=0pt]\item $0\ \mathrm{m/s}$\item $4.8\ \mathrm{m/s}$\item $5.0\ \mathrm{m/s}$\item $10.0\ \mathrm{m/s}$\end{enumerate}
\medskip

\section*{ข้อ 27 \hfill \normalsize MEC-0027}
(สอวน ม.4 2551) อนุภาคเคลื่อนที่ไปตามเส้นรอบวงของวงกลมวงหนึ่งได้หนึ่งในสี่รอบด้วยอัตราเร็วเฉลี่ย $v$ ในช่วงเวลาหนึ่ง จงหาขนาดของความเร็วเฉลี่ยของอนุภาคในช่วงเวลาเดียวกันนี้
\medskip

\section*{ข้อ 28 \hfill \normalsize MEC-0028}
(สอวน ม.5 2552) เรือแล่นเข้าหาหน้าผาด้วยความเร็ว $v$ แล้วเปิดหวูดสั้น ๆ และได้ยินเสียงสะท้อนเมื่อเวลาผ่านไป $t$ หลังเปิดหวูด เรืออยู่ห่างหน้าผาเป็นระยะทางเท่าไรขณะเปิดหวูด (ให้ $c$ เป็นอัตราเร็วของเสียงในอากาศนิ่ง)
\begin{center}\begin{tikzpicture}[>=Stealth,scale=.8]\draw[thick,pattern=north east lines] (0,0) rectangle (7,-.12);\draw[thick] (7,0)--(7,1.5);\fill[gray!25] (7,1.25) rectangle (8,1.4);\node[above] at (7.35,.55){หน้าผา};\node[above] at (5.8,.1){น้ำ\,นิ่ง};\draw[thick,fill=white] (1.4,0)--(2.4,0)--(2.6,.32)--(1.45,.23)--cycle;\draw[->,thick] (2.2,.6)--(3.5,.6) node[midway,above]{$v$};\end{tikzpicture}\end{center}
\medskip

\section*{ข้อ 29 \hfill \normalsize MEC-0029}
(สอวน รอบที่ 1 2547) เรือเคลื่อนที่เข้าหาหน้าผาด้วยความเร็วคงที่ $v$ และส่งเสียงหวูดออกไปขณะที่อยู่ห่างจากหน้าผาเป็นระยะทาง $x$ คนบนเรือต้องรออีกนานเท่าใดหลังส่งเสียงหวูด จึงจะได้ยินเสียงหวูดที่สะท้อนกลับมาจากหน้าผา กำหนดว่าเสียงเคลื่อนที่เร็ว $c$ ในอากาศนิ่ง
\begin{center}\begin{tikzpicture}[>=Stealth,scale=.8]\fill[pattern=north east lines] (0,0) rectangle (.18,1.5);\draw[thick] (.2,.2)--(8,.2);\node[above] at (.65,.9){หน้าผา};\draw[thick,fill=white] (6,.2)--(7,.2)--(7,.52)--(6.1,.6)--cycle;\draw[->,thick] (6.3,.85)--(5,.85) node[midway,above]{$v$};\end{tikzpicture}\end{center}
\medskip

\section*{ข้อ 30 \hfill \normalsize MEC-0030}
(IJSO รอบที่ 2 2554) เรือแล่นเข้าหาหน้าผาด้วยความเร็วคงที่ $u$ เรือเปิดหวูดสั้น ๆ ขณะที่อยู่ห่างจากหน้าผา $L$ อีกนานเท่าใดหลังหวูดนี้คนบนเรือจึงจะได้ยินเสียงสะท้อนครั้งแรก (ให้ใช้ $c$ แทนอัตราเร็วของเสียงในอากาศ)
\begin{center}\begin{tikzpicture}[>=Stealth,scale=.85]\draw[thick] (0,0)--(9,0);\draw[fill=gray!30] (8.4,0) rectangle (8.55,2);\node[left] at (8.4,1.5){หน้าผาชัน};\draw[thick,fill=white] (.6,0)--(2,0)--(2.3,.5)--(.65,.35)--cycle;\node[above] at (1.25,.45){เรือ};\draw[->,thick] (2.5,.9)--(4,.9) node[midway,above]{$u$};\draw[<->] (2, -.4)--(8.4,-.4) node[midway,below]{$L$};\end{tikzpicture}\end{center}
\medskip

\end{document}